\chapter{How the settings are organized}
\label{ch:t-layers}

\section{Five places}

\begin{center}
\begin{tabularx}{\linewidth}{@{}>{\bfseries}p{0.24\linewidth}p{0.3\linewidth}X@{}}
  \toprule
  Place & File & What goes there \\
  \midrule
  Robot profile & \file{x30_bringup/config/robot.yaml} & the body and its motion limits, read by every
    node under \code{robot/}; read-only at runtime \\
  Node settings & \file{x30_<package>/config/*.yaml} & each node's own parameters
    (Appendix~\ref{ch:t-reference}) \\
  Launch arguments & \file{/etc/x30/<role>.env}, \code{X30_LAUNCH_ARGS} & which parts run on a host
    and with what: the site, the controller, NETRA, the camera hosts \\
  Secrets & \file{/etc/x30/<role>.env} & tokens and camera passwords, as environment variables only
    (never parameters: \cmd{roslaunch} prints parameters) \\
  Sites & \file{~/x30/map/<site>/} & routes, stop points, tags and inspection points of a
    substation: taught, not tuned \\
  \bottomrule
\end{tabularx}
\end{center}

The profile is the one place for numbers several nodes must agree on: the footprint, the
speeds, the braking. A node never keeps its own copy of them; where a node needs a stricter value
(the follower's cruise speed below \code{max_vx}), its own parameter can only lower it.

\section{Live and start-only parameters}

Every parameter is declared in a configuration dataclass with its unit, range and description,
and validated when the node starts: a value outside its range or a contradiction (an approach
speed above the cruise speed) stops the node with the reason. The tables of
Appendix~\ref{ch:t-reference} come from the same declarations.

\begin{itemize}
  \item A \textbf{live} parameter is read again every second. Changing it on the parameter server
    takes effect within that second; an invalid value is rejected with a log line and the old
    value stays.
  \item A parameter marked $^\dagger$ is read at start only: change it in the YAML file and
    restart the host's service.
\end{itemize}

\begin{lstlisting}[language=bash]
rosparam set /control_mux_node/field/margin 0.2
rosparam get /control_mux_node/field
\end{lstlisting}

The first is live, in effect within 1~s; the second shows what the node runs with.

\begin{warnbox}[A live change does not survive a restart]
The parameter server is filled from the YAML files at every start. Try a value live, then write
it into the YAML file of the workspace (or the package's) so the next start keeps it.
\end{warnbox}

\section{Where to start}

Tune from the bottom up: each layer relies on the one below being right.

\begin{enumerate}
  \item \textbf{The profile}: measure the base height, the braking and the reaction time
    (Chapter~\ref{ch:t-profile}). Every safety distance follows from them.
  \item \textbf{The driver}: the signs of the axes and the gait's speed limits
    (Chapter~\ref{ch:t-driver}).
  \item \textbf{The scans}: the LiDAR's mount and the height band; a scan that misses an obstacle
    or sees the floor makes every layer above wrong (Chapter~\ref{ch:t-perception}).
  \item \textbf{The protective field}: its margins on the measured braking
    (Chapter~\ref{ch:t-control}).
  \item \textbf{Odometry and localization}: the drift and the tags (Chapters~\ref{ch:t-odometry}
    and~\ref{ch:t-localization}).
  \item \textbf{Route following} and \textbf{avoidance} last (Chapters~\ref{ch:t-following}
    and~\ref{ch:t-obstacles}).
\end{enumerate}

\begin{tipbox}[Record before and after]
A change is judged on data, not on a feeling: \cmd{rosbag record} the topics of the layer before
and after the change, on the same route, and compare. The trip meter
(\topic{odometry/trip}) and \topic{navigation/status} give the time per lap and the stops.
\end{tipbox}
